% ----------------------------------------------------------------------
%  Введение к главе 1
% ----------------------------------------------------------------------
Кристаллография~— язык, на котором говорит всё материаловедение. Чтобы
описать дислокацию (гл.~2), зону Бриллюэна или анизотропию свойств
(гл.~3), текстуру прокатанного листа или ориентационное соотношение при
мартенситном превращении, нужно уметь задавать направления и плоскости в
кристалле, вычислять углы и расстояния между ними, изображать их на
проекциях и понимать, какие ограничения накладывает на свойства
симметрия. Современная кристаллография включает теорию симметрии,
методы определения структуры, изучение реальной структуры, роста и
физических свойств кристаллов. Задача этой главы~— научить
\emph{кристаллографическому подходу} к анализу эксперимента.

\begin{figure}[ht]
\centering
\begin{tikzpicture}[
  box/.style={draw=#1, fill=#1!8, rounded corners=2pt, align=center,
    font=\small, inner sep=4pt, text width=38mm, minimum height=11mm, line width=0.6pt},
  flow/.style={-{Stealth[length=2mm]}, line width=0.6pt, fmGray}]
  \node[box=fmBlue] (a) at (0,0) {\textbf{1.1} Кристаллическое\\состояние};
  \node[box=fmBlue] (b) at (0,-1.6) {\textbf{1.2} Основы\\кристаллографии};
  \node[box=fmGreen] (c) at (-4.6,-3.6) {\textbf{1.3} Симметрия\\кристаллов};
  \node[box=fmGreen] (d) at (-4.6,-5.3) {\textbf{1.4} Структура\\кристаллов};
  \node[box=fmAmber] (e) at (0,-3.6) {\textbf{1.5} Физические\\свойства кристаллов};
  \node[box=fmRed] (f) at (4.6,-3.6) {\textbf{1.6} Кристаллография\\пластической деформации};
  \node[box=fmRed] (g) at (4.6,-5.3) {\textbf{1.7} Границы зёрен};
  \node[box=fmRed] (h) at (4.6,-7.0) {\textbf{1.8} Мартенситные\\превращения};
  \draw[flow] (a) -- (b);
  \draw[flow] (b.west) -| (c.north); \draw[flow] (c) -- (d);
  \draw[flow] (b) -- (e); \draw[flow] (c.east) -- (e.west);
  \draw[flow] (b.east) -| (f.north); \draw[flow] (f) -- (g); \draw[flow] (g) -- (h);
\end{tikzpicture}
\caption{Карта главы. Аппарат кристаллографии (синий блок) нужен для
описания симметрии и структуры (зелёный), физических свойств (жёлтый) и
реальных процессов в металлах~— деформации, границ, превращений
(красный)}
\label{fig:ch1-map}
\end{figure}

\begin{vnimanie}[Как читать эту главу]
Глава насыщена геометрией и матрицами. Все матричные формулы
дублируются простыми частными случаями для кубической решётки~— на
первом проходе можно ограничиться ими. Выводы, нужные лишь для
вычислений «в общем виде» (косоугольные системы, метрические матрицы),
вынесены в блоки «Подробнее».
\end{vnimanie}
